\chapter{Background} 
\label{chap:Background}

AR wearables introduce unique privacy challenges for both the wearer and bystanders. This chapter explains how AR devices operate and the network communication protocols involved, establishing the concepts needed to understand the privacy implications of the Ray-Ban Meta glasses. It also introduces the scope of this thesis.

Augmented reality (AR) combines digital information with the real-time user environment, meaning that it provides a real-world environment with generated perceptual information overlaid on top of it. AR wearables are becoming increasingly sophisticated; for instance, the Ray-Ban Meta glasses have built-in cameras, microphones, and lows for hands-free interactions and real-time analytics. However, these components also allow for continuous data capture from the wearer's perspective, raising concerns about how and when data is being collected, especially for those individuals who are unaware and, hence, do not have the ability to give consent to being recorded.

While traditional Internet of Things (IoT) devices may gather metrics in fixed environments, the Ray-Ban Meta glasses and other AR wearables introduce unique challenges to private and communal privacy. The ability of the Ray-Ban Meta to continuously collect data from the wearer's perspective as they are moving and forward that information to its servers or external services, along with a lack of clarity about the information being gathered, raises ethical concerns about user privacy. 

Moreover, most AR wearables, including the Ray-Ban Meta glasses, rely on a companion app for functionalities, including initial set-up, authentication, firmware updates, and communication with cloud providers or other services. The Meta View app, for instance, can pair with the glasses, import photos or video recordings from the glasses, manage preferences, and potentially share that data with its servers, which leads to a violation of confidentiality as data might be used for analytics beyond the user's awareness. Many devices also rely on external servers, which raises additional concerns, as they may process and retain data for extended periods, making it unclear who has access to user information in the long term or whether it is shared with advertisers. Whether or not the information is shared with advertisers or only with internal analysts, they expand the attack surface of users to a broader network.  

When compromised and intercepted, these communications could expose personal information about the user or bystanders to adversaries. Many AR devices use encryption protocols, such as Transport Layer Security (TLS) and Quick UDP Internet Connections (QUIC), to protect the payload content in transit. However, the accompanying metadata, including packet sizes, timing, and connection endpoints, can still offer insights into user behavior and application usages, allowing attackers to deduce application usage when the user is transmitting data and infer user activities. A blog post, \citet{CodeByter2023}, demonstrated vulnerabilities in the previous generation of the glasses, the Ray-Ban Stories, which allowed interception and decryption of transmitted data. Since then, Meta has implemented improvements to fix those vulnerabilities and secure the communication channels. The blog illustrates how, even with encryption, it is possible to disclose information due to vulnerable or insecure communication channels.

In addition to network traffic, the Ray-Ban Meta glasses use Bluetooth Low Energy (BLE) and Classic Bluetooth for pairing and peripheral management with the companion app. While BLE provides functionalities, attackers monitoring the local environment can capture BLE advertisements, which can expose device names, MAC addresses, and other identifiers, which could link specific wearables to an individual, correlating activities that enhance user tracking. Bluetooth traffic may also contain partial handshakes that inadvertently reveal device capabilities or usage patterns.

Privacy implications extend beyond the wearer but also to those who inadvertently have their images or details collected by the glasses. While the wearer may have consciously decided to use the glasses, bystanders are rarely notified that audio or video recordings are being made. From a bystander's perspective, continuous recording in public spaces confronts conventional notions of informed consent. In addition, bystanders may be subjected to their image or information being sent or shared to cloud providers or advertisers for purposes, including analytics, without any knowledge. The data being sent or shared could also be used in machine learning classifiers or face recognition engines, thus creating more privacy concerns, such as user profiling. Therefore, it is important to have more explicit user and bystander controls, along with stronger privacy protections for both users and bystanders.

This thesis investigates the potential privacy implications of these factors on the Ray-Ban Meta on users' glasses and bystanders by analyzing the device's hardware design, companion app architecture, third-party integrations, and communication protocols. The study covers the static analysis of the app, exploring permissions and code structures that may indicate data flow to external services, network traffic examinations, identifying what is transmitted under different usage scenarios, Bluetooth communications, evaluating device pairing and peripheral interactions for possible vulnerabilities, and implications for bystanders, addressing the privacy concerns for individuals who may not have consented to AR data collection.



%
%\section{IoT}
%
%In addition to network traffic, the glasses rely on Bluetooth communication, using classic Bluetooth and Bluetooth Low Energy (BLE) for different functionalities, such as audio streaming and device pairing.
%
%\section{BLE Pairing}
%In BLE communication, there are specific roles for which devices take on. The first is a central device that scans and initiates connections, which is the role of the phone. The second is a peripheral device that broadcasts their availability for connections, which is the role of the glasses. As a result, they send different packet types when communicating, with the glasses sending packets that announce that it is ready to establish a connection and the phone passively advertising to announce its presence to nearby devices without actively seeking a connection, being in a standby state until pairing is initiated. The amount of traffic kept minimal by having only the devices that are communicating with what is required for their respective roles reflects the energy efficiency design of BLE.
%
%
%\section{Classic Bluetooth Protocol}
%BT BR/EDR RF is a protocol that means that the traffic observed is part of Classic Bluetooth communication, more specifically it refers to the baseband radio frequency layer used in Classical Bluetooth Communication. BR is the Basic Rate of the Classic Bluetooth data that is running at 1 Mbps. EDR is an enhanced BR that supports faster data transfer rates, usually 2 or 3 Mbps). RF refers to Radio Frequency, which is the physical layer of Bluetooth communication that deals with data over a specific frequency band, usually between 2.402 and 2.48 GHz. 

